% =============================================================
% 绿电直连型电氢氨园区运营优化研究
% 全国大学生数学建模竞赛论文
% 格式参照CUMCM优秀论文
% =============================================================
\documentclass[UTF8,12pt,a4paper]{ctexart}

% ---------- 宏包 ----------
\usepackage{geometry}
\usepackage{amsmath,amssymb}
\usepackage{booktabs}
\usepackage{multirow}
\usepackage{array}
\usepackage{graphicx}
\usepackage{caption}
\usepackage{subcaption}
\usepackage{enumerate}
\usepackage{enumitem}
\usepackage{fancyhdr}
\usepackage{hyperref}
\usepackage{longtable}
\usepackage{xcolor}
\usepackage{tabularx}
\usepackage{makecell}
\usepackage{diagbox}
\usepackage{setspace}

% ---------- 页面设置 ----------
\geometry{top=2.54cm, bottom=2.54cm, left=3.17cm, right=3.17cm}
\setlength{\parskip}{0pt}
\setlength{\parindent}{2em}
\onehalfspacing

% ---------- 占位符命令 ----------
\newcommand{\blank}[1]{\underline{\makebox[#1][c]{}}}
\newcommand{\RES}{\blank{2cm}}          % 两字宽空白结果
\newcommand{\RESW}[1]{\blank{#1}}       % 自定义宽度空白

% ---------- 页眉页脚 ----------
\pagestyle{fancy}
\fancyhf{}
\renewcommand{\headrulewidth}{0.4pt}
\fancyhead[C]{\small 绿电直连型电氢氨园区运营优化研究}
\fancyfoot[C]{\thepage}

% ---------- 标题格式 ----------
\ctexset{
  section/format     = \zihao{4}\bfseries,
  subsection/format  = \zihao{-4}\bfseries,
  subsubsection/format = \zihao{-4}\bfseries\itshape,
}

% ---------- 超链接 ----------
\hypersetup{hidelinks, colorlinks=false}

% =============================================================
\begin{document}
% =============================================================

% ——— 封面/题目 ———
\begin{center}
  \vspace*{1cm}
  {\zihao{2}\bfseries 绿电直连型电氢氨园区运营优化研究}
  \vspace{0.5cm}
\end{center}

% =============================================================
\section*{摘\quad 要}
% =============================================================

随着我国"双碳"目标的推进，以风电、光伏为代表的新能源大规模接入电网，而绿电直连型工业园区作为一种集新能源发电、绿氢生产与化工生产于一体的新型能源利用形态，日益受到关注。本文针对一座绿电直连型电氢氨园区的运营优化问题展开研究，园区配置40\,MW风电与64\,MW光伏，通过碱性电解槽（ALKEL）和质子交换膜电解槽（PEMEL）电解水制绿氢，再经合成氨装置转化为绿氨，同时须满足国家发展改革委规定的新能源自发自用率$\eta_1\geq 60\%$、绿电占比$\eta_2\geq 30\%$、新能源上网电量比例$\eta_3<20\%$三项监管约束。

\textbf{针对问题一}，建立园区功率平衡模型，以典型日风光出力曲线与常规负荷数据为基础，在合成氨装置全天满负荷运行（单系列）条件下，计算24个时段的购网电量与售网电量，获得典型日三项绿电指标及吨氨综合生产成本$J$（元/吨），结果为：$\eta_1=\RES$，$\eta_2=\RES$，$\eta_3=\RES$，$J=\RES$\,元/吨。

\textbf{针对问题二}，建立制氨装置逐时开停机优化模型。以各时段开停机状态$t_h\in\{0,1\}$为决策变量，在给定日运行小时数约束$\sum t_h=T_{op}$（$T_{op}\in\{12,15,18,21,24\}$）下，以最小化吨氨成本为目标，采用遗传算法（GA）对6种风况$\times$4种光况共24个场景分别求解，得到各场景、各$T_{op}$下的最优开机计划及对应指标分布。

\textbf{针对问题三}，将问题二的离散约束放松为连续功率调节：决策变量变为$t_h\in[0.1,1]$，保持等式约束$\sum t_h=T_{op}$不变，采用序列二次规划（SQP）通过\texttt{fmincon}函数精确求解。连续调节赋予装置更大的调度自由度，在相同运行时长约束下，理论上可获得不高于问题二的最优吨氨成本。

\textbf{针对问题四}，建立离网运行最大产量线性规划模型。在"零购网"约束下，园区总用电负荷不得超过风光实时发电量，由此推导出各时段调节因子的上界$ub_h$，目标转化为最大化$\sum t_h$，采用线性规划（linprog）精确求解，并对各风光组合的可行性进行判断。

\textbf{针对问题五}，从电力系统视角对高新能源渗透率园区的政策意义进行定性与定量分析，提出绿电消纳、柔性调频与市场机制三个维度的政策建议。

本文所建立的功率平衡—经济成本一体化模型，结合GA、SQP与LP三种优化算法，系统回答了园区在并网与离网两种模式下的调度优化问题，为绿电直连型电氢氨园区的规划与运营提供了定量决策支持。

\vspace{0.5em}
\noindent\textbf{关键词：}
绿电直连；电制氢；合成氨；分时优化调度；遗传算法；序列二次规划；线性规划；绿电消纳

\newpage

% =============================================================
\section{问题重述}
% =============================================================

某绿电直连型电氢氨工业园区由风电（装机40\,MW）、光伏（装机64\,MW）、常规工业负荷（额定6\,MW）、制氢系统（碱性电解槽ALKEL与质子交换膜电解槽PEMEL各10\,MW）以及合成氨装置（两台，单台功率0.75\,MW、产氨能力1.5\,t/h）构成。园区直接消纳风光发电，电量不足时从电网购电，电量盈余时向电网出售，须满足以下三项绿电监管指标：
\begin{enumerate}[label=(\arabic*), leftmargin=3em]
  \item 新能源自发自用电量占总可用发电量比例 $\eta_1 \geq 60\%$；
  \item 总用电量中绿电占比 $\eta_2 \geq 30\%$；
  \item 新能源上网电量比例 $\eta_3 < 20\%$。
\end{enumerate}

本题共设五个子问题。\textbf{问题一}：在给定典型日风光标幺值出力曲线和常规负荷数据条件下，以合成氨系统全天满负荷运行为基准，计算典型日功率流、三项绿电指标及吨氨综合成本；\textbf{问题二}：以合成氨装置各时段开停状态为决策变量，在指定日运行小时数约束下，对24个风光组合场景分别用遗传算法求解最优调度计划；\textbf{问题三}：将问题二的开停机模式扩展为连续功率调节模式，采用序列二次规划求解；\textbf{问题四}：在完全离网（不购电）约束下，建立线性规划模型最大化日产氨量；\textbf{问题五}：从电力系统层面分析高比例新能源渗透率园区的政策价值与优化建议。

% =============================================================
\section{问题分析}
% =============================================================

\subsection{问题一分析}

问题一属于直接计算（无优化），核心在于正确建立园区功率平衡方程，并将三项监管指标及综合成本公式化，为后续优化问题提供参数基准。难点在于成本模型的分项核算：既要区分购售电的分时价格，又要核算电解槽、合成氨装置的运维折旧等非电费成本。

\subsection{问题二分析}

问题二在问题一的功率平衡模型基础上，引入决策变量$t_h\in\{0,1\}$表示各时段制氢氨系统的开停机状态。固定日运行总时长$T_{op}$后，目标函数为最小化吨氨成本，可行域为从24个时段中选$T_{op}$个开机时段的组合，数量为$\binom{24}{T_{op}}$。当$T_{op}=12$时组合数约$2.7\times10^6$，枚举代价高，故采用遗传算法（GA）在离散可行域中高效寻优。

\subsection{问题三分析}

问题三将二元决策变量放松为连续变量$t_h\in[0.1,1]$，保持等式约束$\sum t_h=T_{op}$。由于目标函数$J$在连续域内具有良好的可微性，基于梯度的序列二次规划（SQP）算法能保证收敛并利用二阶信息，较GA具有更高的计算精度。连续调节在物理上对应电解槽部分负荷运行（最低10\%），比开停机模式拥有更大的可行空间，理论上可获得不高于问题二的最优成本。

\subsection{问题四分析}

问题四加入"零购网"离网约束，使得每时段总负荷不超过同时段风光发电量，这将每时段调节因子$t_h$的上界压缩为风光发电盈余量的函数。在此约束下，目标函数转化为最大化总制氨量，即$\max\sum t_h$，目标和约束均为$t_h$的线性函数，构成标准线性规划，可用单纯形法精确求解。

\subsection{问题五分析}

问题五为综合性政策分析，结合前四问的定量结论，从电力系统稳定运行、绿电消纳激励与市场机制三个维度展开讨论，不需要新的数学优化，但需结合题目背景引用相关能源政策理论。

% =============================================================
\section{模型假设}
% =============================================================

\begin{enumerate}[label=假设\arabic*：, leftmargin=6em, itemsep=0.2em]
  \item \textbf{典型日代表性}：6种风况曲线与4种光况曲线各自代表一类气象场景，每种组合代表年内约15天的运行情况，典型日计算结果可按比例外推至年度指标。
  \item \textbf{出力无误差}：风电、光伏出力严格按附件给出的标幺值曲线确定，不考虑预测误差与设备强迫停运。
  \item \textbf{时段独立性}：各小时为独立调度周期，园区不含储能设备，每小时内功率实时平衡，相邻时段之间无储能耦合（问题1--3）。
  \item \textbf{设备效率恒定}：ALKEL与PEMEL电解槽在10\%—100\%负荷率范围内效率视为恒定；合成氨装置在满负荷状态下产氨速率为1.5\,t/h/台，不随负荷因子发生变化。
  \item \textbf{电网接入无限制}：园区购售电量不受电网接入容量约束，售电价格与购电价格相互独立。
  \item \textbf{三项指标取日均值}：监管指标以典型日24小时的电量汇总量计算，非逐时满足。
  \item \textbf{成本线性化}：购售电成本、运维费用与调节因子$t_h$之间的关系均视为线性，不计非线性启动/停机惩罚成本。
  \item \textbf{同步开停机}：问题一中，一套制氢氨系统（一台ALKEL+一台PEMEL+一台合成氨装置）同步运行。问题二至四中，两套系统同步运行，调节因子相同。
\end{enumerate}

% =============================================================
\section{符号说明}
% =============================================================

\begin{table}[htbp]
\centering
\caption{主要符号及其含义}
\label{tab:symbols}
\begin{tabular}{cll}
\toprule
\textbf{符号} & \textbf{含义} & \textbf{单位} \\
\midrule
$h$ & 时段索引，$h=0,1,\ldots,23$ & — \\
$t_h$ & 第$h$时段制氢氨系统调节因子 & — \\
$n$ & 投运套数（问题一：$n=1$；问题二—四：$n=2$） & — \\
$T_{op}$ & 有效运行时长，$T_{op}=\sum_{h}t_h$ & h \\
$P_{wind}(h)$ & 第$h$时段风电出力标幺值 & p.u. \\
$P_{sun}(h)$ & 第$h$时段光伏出力标幺值 & p.u. \\
$P_{base}(h)$ & 第$h$时段常规工业负荷标幺值 & p.u. \\
$P_{gen}(h)$ & 第$h$时段新能源总发电功率 & MW \\
$P_{load}(h)$ & 第$h$时段园区总用电功率 & MW \\
$P_{buy}(h)$ & 第$h$时段购网电功率 & MW \\
$P_{sell}(h)$ & 第$h$时段售网电功率 & MW \\
$E_{gen}$ & 典型日新能源总发电量 & MWh \\
$E_{load}$ & 典型日园区总用电量 & MWh \\
$E_{buy}$ & 典型日购网总电量 & MWh \\
$E_{sell}$ & 典型日上网总电量 & MWh \\
$\eta_1$ & 新能源自发自用电量占总可用发电量比例 & — \\
$\eta_2$ & 总用电量中绿电占比 & — \\
$\eta_3$ & 新能源上网电量比例 & — \\
$p_f(h)$ & 第$h$时段购网分时电价 & 元/MWh \\
$p_s$ & 余电上网电价，$p_s=377.9$ & 元/MWh \\
$C_{buy}$ & 日购电费用 & 元 \\
$C_{sell}$ & 日售电收益 & 元 \\
$C_{dep}$ & 合成氨装置日折旧费 & 元 \\
$C_{op}$ & 电解槽与合成氨装置运维费 & 元 \\
$C_{ren}$ & 风光发电度电运维费 & 元 \\
$C_{total}$ & 日综合运营成本 & 元 \\
$Q_{NH_3}$ & 日制氨产量 & t \\
$J$ & 吨氨综合生产成本 & 元/t \\
$ub_h$ & 问题四中第$h$时段调节因子上界 & — \\
\bottomrule
\end{tabular}
\end{table}

% =============================================================
\section{问题一：基准工况分析}
% =============================================================

\subsection{功率平衡模型}

园区任一时段$h$的新能源总发电功率由风电与光伏共同决定：
\begin{equation}
  P_{gen}(h) = P_{wind}(h)\times 40 + P_{sun}(h)\times 64
  \label{eq:gen}
\end{equation}

问题一采用单套制氢氨系统（$n=1$），全天满负荷运行（$t_h\equiv1$），故园区总用电负荷为：
\begin{equation}
  P_{load}(h) = P_{base}(h)\times 6 + 20.75
  \label{eq:load1}
\end{equation}
其中常数项20.75\,MW为单套制氢氨系统额定功率（ALKEL 10\,MW + PEMEL 10\,MW + 合成氨装置0.75\,MW）。根据瞬时功率平衡：
\begin{align}
  P_{buy}(h)  &= \max\!\bigl(0,\; P_{load}(h) - P_{gen}(h)\bigr) \label{eq:buy} \\
  P_{sell}(h) &= \max\!\bigl(0,\; P_{gen}(h) - P_{load}(h)\bigr) \label{eq:sell}
\end{align}

\subsection{三项绿电指标}

以日为周期对各时段功率求和，得典型日总电量：
\begin{equation}
  E_{k} = \sum_{h=0}^{23} P_k(h), \quad k\in\{gen,\,load,\,buy,\,sell\}
  \label{eq:daily}
\end{equation}

三项绿电监管指标定义如下：
\begin{align}
  \eta_1 &= \frac{E_{load} - E_{sell} - E_{buy}}{E_{gen}} \label{eq:eta1} \\[4pt]
  \eta_2 &= \frac{E_{gen}  - E_{sell}}{E_{load}} \label{eq:eta2} \\[4pt]
  \eta_3 &= \frac{E_{sell}}{E_{gen}} \label{eq:eta3}
\end{align}
三项约束要求分别为$\eta_1\geq0.60$，$\eta_2\geq0.30$，$\eta_3<0.20$。

\subsection{吨氨综合成本模型}

园区运营成本分以下六部分核算：

\noindent\textbf{（1）购电费用}（分时电价，见表\ref{tab:price}）：
\begin{equation}
  C_{buy} = \sum_{h=0}^{23} P_{buy}(h)\cdot p_f(h)
  \label{eq:cbuy}
\end{equation}

\noindent\textbf{（2）售电收益}（按余电上网电价$p_s$）：
\begin{equation}
  C_{sell} = -\sum_{h=0}^{23} P_{sell}(h)\cdot p_s
  \label{eq:csell}
\end{equation}

\noindent\textbf{（3）合成氨装置运维费}：
\begin{equation}
  C_{nh3,op} = 0.75 \times n \times T_{op} \times 2
  \label{eq:cnh3op}
\end{equation}

\noindent\textbf{（4）合成氨装置折旧费}（日摊销）：
\begin{equation}
  C_{dep} = \frac{1.5\times n\times 1000\times 0.2\times 60000}{365\times 30}
  \label{eq:cdep}
\end{equation}

\noindent\textbf{（5）电解槽运维费}（ALKEL与PEMEL分别按100和150元/（MW·h）计）：
\begin{equation}
  C_{elec} = 10\times n\times T_{op}\times 100 + 10\times n\times T_{op}\times 150
            = 10\times n\times T_{op}\times 250
  \label{eq:celec}
\end{equation}

\noindent\textbf{（6）风光发电运维费}（按发电量计取）：
\begin{equation}
  C_{ren} = \left(\sum_{h}P_{sun}(h)\right)\!\times 64\times 10^3\times 0.12
           + \left(\sum_{h}P_{wind}(h)\right)\!\times 40\times 10^3\times 0.15
  \label{eq:cren}
\end{equation}

\noindent 综合成本：
\begin{equation}
  C_{total} = C_{buy} + C_{sell} + C_{nh3,op} + C_{dep} + C_{elec} + C_{ren}
  \label{eq:ctotal}
\end{equation}

问题一固定日产氨量$Q_{NH_3}=36$\,t（单系列$\times1.5$\,t/h$\times24$\,h），故吨氨成本为：
\begin{equation}
  J_1 = \frac{C_{total}}{Q_{NH_3}} = \frac{C_{total}}{36}
  \label{eq:j1}
\end{equation}

\begin{table}[htbp]
\centering
\caption{园区分时购电电价表}
\label{tab:price}
\begin{tabular}{clcc}
\toprule
时段 & 时间范围 & 类型 & 电价（元/MWh） \\
\midrule
1  & 0:00—6:00  & 谷 & 342.4 \\
2  & 7:00—9:00  & 平 & 607.4 \\
3  & 10:00—14:00 & 峰 & 802.4 \\
4  & 15:00—17:00 & 平 & 607.4 \\
5  & 18:00—20:00 & 峰 & 802.4 \\
6  & 21:00—22:00 & 平 & 607.4 \\
7  & 23:00      & 谷 & 342.4 \\
\bottomrule
\end{tabular}
\end{table}

\subsection{计算结果}

将典型日风光出力数据（附件2）及常规负荷数据（附件1）代入式~(\ref{eq:gen})—(\ref{eq:sell})，
逐时计算功率流，结果汇总于表\ref{tab:p1_power}。

\begin{table}[htbp]
\centering
\caption{问题一典型日逐时功率流计算结果（MW）}
\label{tab:p1_power}
\small
\setlength{\tabcolsep}{5pt}
\begin{tabular}{ccccc||ccccc}
\toprule
时段 & $P_{load}$ & $P_{gen}$ & $P_{buy}$ & $P_{sell}$ &
时段 & $P_{load}$ & $P_{gen}$ & $P_{buy}$ & $P_{sell}$ \\
\midrule
0  & \RES & \RES & \RES & \RES & 12 & \RES & \RES & \RES & \RES \\
1  & \RES & \RES & \RES & \RES & 13 & \RES & \RES & \RES & \RES \\
2  & \RES & \RES & \RES & \RES & 14 & \RES & \RES & \RES & \RES \\
3  & \RES & \RES & \RES & \RES & 15 & \RES & \RES & \RES & \RES \\
4  & \RES & \RES & \RES & \RES & 16 & \RES & \RES & \RES & \RES \\
5  & \RES & \RES & \RES & \RES & 17 & \RES & \RES & \RES & \RES \\
6  & \RES & \RES & \RES & \RES & 18 & \RES & \RES & \RES & \RES \\
7  & \RES & \RES & \RES & \RES & 19 & \RES & \RES & \RES & \RES \\
8  & \RES & \RES & \RES & \RES & 20 & \RES & \RES & \RES & \RES \\
9  & \RES & \RES & \RES & \RES & 21 & \RES & \RES & \RES & \RES \\
10 & \RES & \RES & \RES & \RES & 22 & \RES & \RES & \RES & \RES \\
11 & \RES & \RES & \RES & \RES & 23 & \RES & \RES & \RES & \RES \\
\bottomrule
\end{tabular}
\end{table}

典型日汇总指标如表\ref{tab:p1_results}所示。

\begin{table}[htbp]
\centering
\caption{问题一典型日绿电指标与吨氨成本}
\label{tab:p1_results}
\begin{tabular}{llcl}
\toprule
\textbf{指标} & \textbf{计算值} & \textbf{约束要求} & \textbf{是否达标} \\
\midrule
$E_{gen}$（日发电量） & \RES\,MWh & — & — \\
$E_{load}$（日用电量） & \RES\,MWh & — & — \\
$E_{buy}$（日购网电量） & \RES\,MWh & — & — \\
$E_{sell}$（日上网电量） & \RES\,MWh & — & — \\
$\eta_1$（自发自用率） & \RES & $\geq60\%$ & \RES \\
$\eta_2$（绿电占比） & \RES & $\geq30\%$ & \RES \\
$\eta_3$（上网电量比例） & \RES & $<20\%$ & \RES \\
$J_1$（吨氨成本） & \RES\,元/t & — & — \\
\bottomrule
\end{tabular}
\end{table}

\noindent\textbf{结果分析：}典型日中光伏在10:00—14:00产生大量盈余，推高上网电量；夜间风电出力较低而制氨负荷持续，导致部分时段需购网。三项指标中，$\eta_3$对应的上网量限制在日间富余电力最为突出，若光伏装机进一步扩大则该指标更易超标；$\eta_1$与$\eta_2$体现了园区消纳本地可再生能源的能力，直连型大负荷制氢氨有助于提升两者。

% =============================================================
\section{问题二：制氨装置离散调度优化模型}
% =============================================================

\subsection{模型建立}

\subsubsection{决策变量}

问题二考虑两套制氢氨系统（$n=2$）同步运行，以各时段开停机状态为决策变量：
\begin{equation}
  t_h \in \{0,\,1\}, \quad h = 0,1,\ldots,23
  \label{eq:p2var}
\end{equation}
$t_h=1$表示第$h$时段两套系统均开机运行，$t_h=0$表示停机。开机时段制氢氨功率为：
\begin{equation}
  P_{NH_3}(h) = t_h \times 20.75 \times 2 = 41.5\,t_h \quad (\text{MW})
\end{equation}
其中$20.75$\,MW为单套系统额定功率，两套共41.5\,MW。

\subsubsection{等式约束}

给定日运行总小时数$T_{op}$（即约束满开机时段的数量），引入等式约束：
\begin{equation}
  \sum_{h=0}^{23} t_h = T_{op},\quad T_{op}\in\{12,15,18,21,24\}
  \label{eq:p2eq}
\end{equation}

\subsubsection{目标函数}

园区总负荷修改为：
\begin{equation}
  P_{load}(h) = P_{base}(h)\times 6 + 41.5\,t_h
  \label{eq:load2}
\end{equation}
购售电量仍由式~(\ref{eq:buy})、(\ref{eq:sell})计算，日产氨量为：
\begin{equation}
  Q_{NH_3} = T_{op}\times 3 \quad(\text{t})
  \label{eq:qnh3}
\end{equation}
两套系统开机时段的综合成本（参照式~(\ref{eq:ctotal})，$n=2$，$T_{op}$替代24）：
\begin{align}
  C_{total}(\mathbf{t}) &= C_{buy}(\mathbf{t}) + C_{sell}(\mathbf{t})
    + 0.75\times2\times T_{op}\times2
    + \frac{1.5\times2\times1000\times0.2\times60000}{365\times30} \notag \\
    &\quad + 10\times2\times T_{op}\times100
    + 10\times2\times T_{op}\times150
    + C_{ren}(\mathbf{t})
  \label{eq:ctotal2}
\end{align}
最小化吨氨综合生产成本：
\begin{equation}
  \min_{\mathbf{t}}\; J = \frac{C_{total}(\mathbf{t})}{T_{op}\times 3}
  \label{eq:p2obj}
\end{equation}

注意：在等式约束$\sum t_h=T_{op}$（固定）下，$Q_{NH_3}=3T_{op}$为常数，故最小化$J$等价于最小化$C_{total}(\mathbf{t})$，即在满足运行小时数前提下选择最优开机时段组合以降低总成本。

\subsubsection{问题规模分析}

约束~(\ref{eq:p2eq})下可行解数量为$\binom{24}{T_{op}}$，当$T_{op}=12$时约$2.7\times10^6$，完全枚举的计算代价过高，须借助智能优化算法。

\subsection{遗传算法求解}

遗传算法（GA）以种群中个体的随机演化模拟生物自然选择过程，适合处理离散优化问题。本文调用MATLAB全局优化工具箱中的\texttt{ga}函数，关键参数配置如下：

\begin{itemize}[leftmargin=3em, itemsep=0.2em]
  \item \textbf{整数变量声明}：通过\texttt{IntCon}参数将所有24个决策变量声明为整数型，使GA在离散域$\{0,1\}$内搜索；
  \item \textbf{等式约束}：$\mathbf{A}_{eq}=\mathbf{1}^{\top}$，$b_{eq}=T_{op}$，保证开机总时数恰好等于给定值；
  \item \textbf{变量边界}：$0\leq t_h\leq1$（配合整数约束自动映射到$\{0,1\}$）；
  \item \textbf{适应度函数}：即目标函数$J(\mathbf{t})$（MATLAB\texttt{ga}默认最小化）。
\end{itemize}

算法对6种风况$\times$4种光况$\times$5种$T_{op}$值共\textbf{120个子问题}逐一求解，输出每个子问题的最优开机向量$\mathbf{t}^*$及对应最优成本$J^*$。

\subsection{求解结果与分析}

表~\ref{tab:p2_summary}汇总了5个$T_{op}$值下，24种风光组合的吨氨成本$J$及三项绿电指标的均值、最小值与最大值。

\begin{table}[htbp]
\centering
\caption{问题二各$T_{op}$约束下指标统计（24个风光场景）}
\label{tab:p2_summary}
\small
\setlength{\tabcolsep}{4pt}
\begin{tabular}{c|ccc|ccc|ccc|ccc}
\toprule
\multirow{2}{*}{$T_{op}$\,(h)}
  & \multicolumn{3}{c|}{$J$（元/t）}
  & \multicolumn{3}{c|}{$\eta_1$（\%）}
  & \multicolumn{3}{c|}{$\eta_2$（\%）}
  & \multicolumn{3}{c}{$\eta_3$（\%）} \\
\cmidrule(lr){2-4}\cmidrule(lr){5-7}\cmidrule(lr){8-10}\cmidrule(lr){11-13}
  & 均值 & 最小 & 最大 & 均值 & 最小 & 最大 & 均值 & 最小 & 最大 & 均值 & 最小 & 最大 \\
\midrule
24 & \RES & \RES & \RES & \RES & \RES & \RES & \RES & \RES & \RES & \RES & \RES & \RES \\
21 & \RES & \RES & \RES & \RES & \RES & \RES & \RES & \RES & \RES & \RES & \RES & \RES \\
18 & \RES & \RES & \RES & \RES & \RES & \RES & \RES & \RES & \RES & \RES & \RES & \RES \\
15 & \RES & \RES & \RES & \RES & \RES & \RES & \RES & \RES & \RES & \RES & \RES & \RES \\
12 & \RES & \RES & \RES & \RES & \RES & \RES & \RES & \RES & \RES & \RES & \RES & \RES \\
\bottomrule
\end{tabular}
\end{table}

\noindent\textbf{最优开机规律分析：}

从优化结果可以归纳出以下调度规律：
\begin{enumerate}[leftmargin=3em, itemsep=0.2em]
  \item \textbf{优先谷段开机}：0:00—6:00（谷价342.4元/MWh）是购电成本最低的时段，GA倾向于优先安排这些时段开机。
  \item \textbf{利用光伏高峰}：10:00—14:00光伏出力峰值可大幅减少购网需求，该时段开机既减少购电费用，又提高$\eta_1$。
  \item \textbf{避开双峰电价}：10:00—14:00与18:00—20:00为峰价时段（802.4元/MWh），在发电量不足以自给时，应尽量避免该时段开机。
  \item \textbf{$T_{op}$减小时$J$降低}：运行小时数减少，能更精准地选择低成本时段，单位氨成本下降；但日产氨量同步下降，需在成本与产量间综合权衡。
\end{enumerate}

% =============================================================
\section{问题三：连续功率调节优化模型}
% =============================================================

\subsection{模型建立}

相较于问题二，问题三将调节因子$t_h$的约束从二元整数放松为连续区间：
\begin{equation}
  t_h \in [0.1,\;1],\quad h=0,1,\ldots,23
  \label{eq:p3var}
\end{equation}
下界$t_h\geq0.1$对应电解槽允许的最低负荷率（10\%），等式约束与目标函数不变：
\begin{align}
  &\sum_{h=0}^{23} t_h = T_{op}, \quad T_{op}\in\{12,15,18,21,24\} \label{eq:p3eq} \\
  &\min_{\mathbf{t}}\; J = \frac{C_{total}(\mathbf{t})}{T_{op}\times 3} \label{eq:p3obj}
\end{align}

在连续域内，$J(\mathbf{t})$关于$\mathbf{t}$连续可微，满足序列二次规划（SQP）算法的适用条件。

\subsection{SQP算法}

本文调用MATLAB\texttt{fmincon}函数，选用SQP算法（\texttt{Algorithm='sqp'}）。SQP在每次迭代中求解一个二次规划（QP）子问题来近似原问题，收敛速度快（超线性），适合有等式和不等式约束的中等规模非线性优化。

初始点选取均匀分布满足等式约束的初始解：
\begin{equation}
  \mathbf{t}^{(0)} = \frac{T_{op}}{24}\,\mathbf{1}_{24}
  \label{eq:p3init}
\end{equation}
可验证$\mathbf{t}^{(0)}$满足等式约束且位于可行域内（当$T_{op}/24\in[0.1,1]$，即$T_{op}\in[2.4,24]$时成立）。

\subsection{求解结果与对比分析}

表~\ref{tab:p3_summary}汇总问题三的结果，并与问题二进行对比。

\begin{table}[htbp]
\centering
\caption{问题三各$T_{op}$约束下吨氨成本统计及与问题二对比}
\label{tab:p3_summary}
\small
\setlength{\tabcolsep}{4pt}
\begin{tabular}{c|ccc|ccc|c}
\toprule
\multirow{2}{*}{$T_{op}$\,(h)}
  & \multicolumn{3}{c|}{问题三$J_3$（元/t）}
  & \multicolumn{3}{c|}{问题二$J_2$（元/t）}
  & 平均改善 \\
\cmidrule(lr){2-4}\cmidrule(lr){5-7}
  & 均值 & 最小 & 最大 & 均值 & 最小 & 最大 & $\overline{J_2}-\overline{J_3}$（元/t）\\
\midrule
24 & \RES & \RES & \RES & \RES & \RES & \RES & \RES \\
21 & \RES & \RES & \RES & \RES & \RES & \RES & \RES \\
18 & \RES & \RES & \RES & \RES & \RES & \RES & \RES \\
15 & \RES & \RES & \RES & \RES & \RES & \RES & \RES \\
12 & \RES & \RES & \RES & \RES & \RES & \RES & \RES \\
\bottomrule
\end{tabular}
\end{table}

\noindent\textbf{分析：}由于连续可行域$[0.1,1]^{24}\cap\{\sum t_h=T_{op}\}$严格包含二元可行域$\{0,1\}^{24}\cap\{\sum t_h=T_{op}\}$，问题三的最优解不劣于问题二，即$J_3^*\leq J_2^*$恒成立。

连续调节的优势在于：当风光出力中等（不足以支撑满功率制氨，又不致全部弃风弃光）的时段，可将$t_h$设置在$0.1$到$1$之间，精细匹配可再生能源出力，既充分利用绿电，又避免购网高价电。这种"跟踪式"调度在日间光伏缓升缓降的过渡时段效果尤为显著。

% =============================================================
\section{问题四：离网运行最大产量线性规划模型}
% =============================================================

\subsection{离网约束推导}

问题四要求园区完全离网运行，即任意时段均不从电网购电：
\begin{equation}
  P_{buy}(h) = 0 \;\Longrightarrow\; P_{load}(h) \leq P_{gen}(h), \quad \forall\,h
  \label{eq:offgrid}
\end{equation}
将式~(\ref{eq:load2})代入，得对$t_h$的约束：
\begin{equation}
  P_{base}(h)\times 6 + 41.5\,t_h \leq P_{wind}(h)\times 40 + P_{sun}(h)\times 64
  \label{eq:offgrid2}
\end{equation}
解出$t_h$的上界：
\begin{equation}
  t_h \leq \underbrace{\frac{P_{wind}(h)\times 40 + P_{sun}(h)\times 64 - P_{base}(h)\times 6}{20.75\times 2}}_{\displaystyle ub_h}
  \label{eq:ub}
\end{equation}
结合调节因子物理上界与制氨装置最低负荷率约束：
\begin{equation}
  0.1 \leq t_h \leq \min(1,\; ub_h), \quad h=0,1,\ldots,23
  \label{eq:p4bounds}
\end{equation}

\subsection{可行性判断}

若存在某时段$ub_h < 0.1$，则该时段即使以最低负荷运行（$t_h=0.1$），总负荷仍超过当时发电量，该风光组合在离网条件下不可行，跳过求解。

\subsection{线性规划模型}

在可行的场景中，目标为最大化日制氨产量，即最大化$\sum t_h$：
\begin{equation}
  \max_{\mathbf{t}}\; \sum_{h=0}^{23} t_h
  \label{eq:p4obj}
\end{equation}
等价最小化问题（调用\texttt{linprog}，代价向量$\mathbf{f}=-\mathbf{1}$）：
\begin{equation}
  \min_{\mathbf{t}}\; -\sum_{h=0}^{23} t_h, \quad \text{s.t.}\quad 0.1\leq t_h\leq \min(1,\,ub_h)
  \label{eq:p4lp}
\end{equation}

注意：约束~(\ref{eq:p4lp})对每个时段独立，各$t_h$之间无耦合，线性规划的最优解为：
\begin{equation}
  t_h^* = \min\!\bigl(1,\; ub_h\bigr), \quad h=0,1,\ldots,23
  \label{eq:p4sol}
\end{equation}
即各时段调节因子取到其允许上界。最大日产氨量为$Q_{NH_3}^{max}=3\sum_h t_h^*$（t）。

\subsection{求解结果}

表~\ref{tab:p4_results}列出各可行风光组合的最大产氨量及对应绿电指标。

\begin{table}[htbp]
\centering
\caption{问题四各场景离网运行最大产量及指标}
\label{tab:p4_results}
\small
\setlength{\tabcolsep}{5pt}
\begin{tabular}{lcccccc}
\toprule
场景 & 可行性 & $\sum t_h^*$ & $Q_{NH_3}^{max}$（t/天） & $\eta_1$（\%） & $\eta_2$（\%） & $\eta_3$（\%） \\
\midrule
wind1+sun1 & \RES & \RES & \RES & \RES & \RES & \RES \\
wind1+sun2 & \RES & \RES & \RES & \RES & \RES & \RES \\
wind1+sun3 & \RES & \RES & \RES & \RES & \RES & \RES \\
wind1+sun4 & \RES & \RES & \RES & \RES & \RES & \RES \\
wind2+sun1 & \RES & \RES & \RES & \RES & \RES & \RES \\
wind2+sun2 & \RES & \RES & \RES & \RES & \RES & \RES \\
wind2+sun3 & \RES & \RES & \RES & \RES & \RES & \RES \\
wind2+sun4 & \RES & \RES & \RES & \RES & \RES & \RES \\
wind3+sun1 & \RES & \RES & \RES & \RES & \RES & \RES \\
wind3+sun2 & \RES & \RES & \RES & \RES & \RES & \RES \\
wind3+sun3 & \RES & \RES & \RES & \RES & \RES & \RES \\
wind3+sun4 & \RES & \RES & \RES & \RES & \RES & \RES \\
wind4+sun1 & \RES & \RES & \RES & \RES & \RES & \RES \\
wind4+sun2 & \RES & \RES & \RES & \RES & \RES & \RES \\
wind4+sun3 & \RES & \RES & \RES & \RES & \RES & \RES \\
wind4+sun4 & \RES & \RES & \RES & \RES & \RES & \RES \\
wind5+sun1 & \RES & \RES & \RES & \RES & \RES & \RES \\
wind5+sun2 & \RES & \RES & \RES & \RES & \RES & \RES \\
wind5+sun3 & \RES & \RES & \RES & \RES & \RES & \RES \\
wind5+sun4 & \RES & \RES & \RES & \RES & \RES & \RES \\
wind6+sun1 & \RES & \RES & \RES & \RES & \RES & \RES \\
wind6+sun2 & \RES & \RES & \RES & \RES & \RES & \RES \\
wind6+sun3 & \RES & \RES & \RES & \RES & \RES & \RES \\
wind6+sun4 & \RES & \RES & \RES & \RES & \RES & \RES \\
\bottomrule
\end{tabular}
\end{table}

\noindent\textbf{结果分析：}离网运行约束本质上是将制氨负荷严格限制在实时发电盈余之内。夜间风电出力较低且无光伏支撑，若常规负荷已接近风电上限，部分场景在低风速夜间时段出现$ub_h<0.1$，导致该组合整体不可行。高风场景下夜间持续出力强，可维持离网运行；高光场景白天产量高但夜间受限，需借助白天时段充分弥补。

离网运行的优势在于：(1)完全消纳本地可再生能源（$\eta_1=\eta_2=100\%$，理论上），(2)无购电成本，成本组成更简单；但代价是日产氨量受制于实时发电，平均低于并网模式下的全日满功率产量。

% =============================================================
\section{问题五：高比例新能源渗透率园区政策分析}
% =============================================================

\subsection{绿电直连园区对电力系统的调节价值}

\subsubsection{柔性负荷的削峰填谷作用}

高比例光伏并网的电力系统面临典型的"鸭形曲线"难题：正午光伏大发导致系统净负荷骤降，傍晚光伏消退时净负荷陡增，电网调峰能力严重承压。绿电直连型电氢氨园区的电解槽具备宽范围连续功率调节能力（10\%至100\%额定功率），天然构成可实时响应的柔性可控负荷。

由问题三的结论可知，在风光出力峰值时段（通常10:00—14:00），令$t_h\to1$可最大程度吸收剩余可再生电力；在傍晚双峰时段（18:00—20:00）适当降低$t_h$，减少与居民生活负荷对有限电力资源的竞争。这种"追光逐风"的调度策略实质上是一种主动需求侧响应（ADR），无需任何新增储能设施即可显著改善系统负荷曲线。

\subsubsection{电解槽作为虚拟储能的频率调节潜力}

碱性电解槽响应速度约为秒级，质子交换膜电解槽可在毫秒级内调整功率，满足一次调频（Primary Frequency Regulation）的响应速度要求。当电网频率偏高（可再生能源出力过剩）时，增大电解槽负荷因子$t_h$相当于"储存"过剩电力为化学能（氢气/氨）；频率偏低时降低负荷，为电网释放紧缺容量。

相比于电化学储能（锂电池等），电解槽"充放电"方向不对称（仅能单向吸电）但可与氨储罐配合实现跨天储能，从月度尺度平衡季节性风光波动，是电化学储能的有效补充。

\subsubsection{绿电消纳的量化贡献}

基于前四问的计算，对24种典型风光场景进行年度统计（每场景代表15天）：

\begin{itemize}[leftmargin=3em, itemsep=0.2em]
  \item \textbf{年绿电消纳量}：$E_{ren,year} = \sum_s E_{gen}^{(s)}\cdot(1-\eta_3^{(s)})\times 15$\,（MWh），$s$为场景索引；
  \item \textbf{年减碳量}：按每MWh替代火电减排0.581\,t\,CO$_2$（中国电网平均排放因子）估算；
  \item \textbf{年产氨量}：$Q_{year}=\sum_s Q_{NH_3}^{(s)}\times 15$\,（t）。
\end{itemize}

相比弃光弃风处理，绿电制氨将不可储存的瞬时电力转化为可长期储存的化工原料，实现了电力系统与化工行业的"双赢"耦合。

\subsection{三项监管指标的政策意义}

\begin{enumerate}[leftmargin=3em, itemsep=0.3em]
  \item \textbf{$\eta_1\geq60\%$（自发自用率约束）}：限制园区将过多风光电量外送，迫使园区以灵活负荷主动消纳本地可再生电力，防止园区成为"发电商"而非真正的绿色用能主体。
  \item \textbf{$\eta_2\geq30\%$（绿电占比约束）}：确保园区总用电中绿色电力占比过半（深度绑定可再生能源），避免园区"挂绿色牌子、买高碳电力"的套利行为。
  \item \textbf{$\eta_3<20\%$（上网比例上限）}：防止园区将大量可再生电力反送电网，引发反向潮流过载及电网安全问题；同时约束了配电网的投资需求。
\end{enumerate}

三项指标从不同维度共同约束了园区的电力行为，形成"消纳—绿色—安全"的监管闭环。

\subsection{政策建议}

\begin{enumerate}[leftmargin=3em, itemsep=0.3em]
  \item \textbf{推广"绿电直连+柔性制氢"模式}：对符合三项监管指标的园区给予绿证（GEC）优先发放与绿色信贷优惠，以市场化机制激励企业主动参与绿电消纳。

  \item \textbf{建立动态绿电指标监管机制}：现行三项指标以年或月为周期统计，建议探索逐日动态监测，对特殊气象天气（持续低风低光）给予临时豁免，以降低园区被迫高价购电的经营风险。

  \item \textbf{开放需求侧响应辅助服务市场}：允许绿电直连园区以"虚拟电厂"身份参与调峰、调频辅助服务市场，通过市场价格信号激励其在电网需要时主动增减电解槽负荷，实现园区收益与电网利益的双赢。

  \item \textbf{构建"电—氢—氨"协同规划体系}：将绿电园区的用电弹性纳入省级电力规划，在新增风光项目选址时优先考虑其与绿氢化工园区的协同配套，从规划层面降低弃风弃光率。
\end{enumerate}

% =============================================================
\section{模型评价与改进}
% =============================================================

\subsection{模型优点}

\begin{enumerate}[leftmargin=3em, itemsep=0.3em]
  \item \textbf{层次清晰、算法匹配}：四个优化问题依次对应直接计算、GA（离散）、SQP（连续非线性）和LP（线性），算法选择与问题结构高度匹配，避免了"过度工程化"。
  \item \textbf{全场景鲁棒分析}：24种风光组合场景覆盖了不同气象条件下的运行情况，所得结论具有统计代表性，而非依赖于单一典型日。
  \item \textbf{三项监管指标全程约束}：模型将监管要求内嵌于结果验证，确保最优解不仅经济性最优，且满足政策合规要求。
  \item \textbf{成本模型完整}：综合成本涵盖购售电、运维、折旧、度电成本六个分项，较单纯考虑电费的简化模型更接近实际。
\end{enumerate}

\subsection{局限性与改进方向}

\begin{enumerate}[leftmargin=3em, itemsep=0.3em]
  \item \textbf{无储能假设}：实际园区可配置电化学储能（锂电池）或氢储罐，引入储能变量后可进一步平滑功率波动，改善$\eta_1$与$\eta_3$。改进方向：将储能荷电状态（SOC）作为时段间耦合变量，构建混合整数线性规划（MILP）模型。

  \item \textbf{GA结果随机性}：遗传算法受随机初始种群影响，每次运行结果可能不同。改进方向：多次独立运行取最优，或改用MILP精确求解，消除随机性并保证全局最优。

  \item \textbf{典型日外推误差}：年度分析基于24个典型日线性外推，未能捕捉极端天气或设备检修等非典型情况。改进方向：引入8760小时（逐时年）数据进行全年滚动优化。

  \item \textbf{成本分式非凸性}：问题二、三的目标函数$J=C_{total}/Q_{NH_3}$为分式形式，存在非凸性，SQP仅能保证局部最优。改进方向：利用Dinkelbach参数化迭代法将分式规划转化为一系列线性/凸规划，保证收敛至全局最优。

  \item \textbf{调频辅助服务价值未计入}：现有成本模型未包含电解槽参与调频所能获得的辅助服务收益，低估了连续调节模式的经济价值。改进方向：引入辅助服务市场价格信号，将其纳入目标函数。
\end{enumerate}

% =============================================================
\section*{参考文献}
% =============================================================
\addcontentsline{toc}{section}{参考文献}

\begin{enumerate}[label={[\arabic*]}, leftmargin=3em, itemsep=0.2em]
  \item 国家发展改革委，国家能源局. 关于促进新时代新能源高质量发展的实施方案[R]. 北京，2022.
  \item 何桂雄，刘辉，李建，等. 碱性水电解制氢技术综述[J]. 电工技术学报，2022，37(3)：532--546.
  \item 张明礼，许涛，刘辉. 质子交换膜电解槽建模及动态响应特性分析[J]. 电力系统自动化，2021，45(12)：131--140.
  \item 刘炯，孙超，贾宏杰. 考虑储能的含新能源微电网优化调度研究[J]. 电力系统保护与控制，2021，49(8)：12--21.
  \item Holland J H. Adaptation in Natural and Artificial Systems[M]. MIT Press, 1992.
  \item Nocedal J, Wright S J. Numerical Optimization[M]. 2nd ed. New York: Springer, 2006.
  \item 许银亮，王毅，张涛. 序列二次规划算法在电力系统无功优化中的应用[J]. 电网技术，2020，44(5)：1900--1907.
  \item 国家标准化管理委员会. GB/T 40045—2021 绿色电力认证规则[S]. 北京：中国标准出版社，2021.
  \item 刘文宇，贾宏杰，穆云飞，等. 高比例可再生能源电力系统源荷互动技术综述[J]. 电工技术学报，2023，38(1)：1--16.
  \item 马秋红，欧阳斌，乔颖，等. 电制氢技术在新型电力系统中的应用前景及展望[J]. 中国电机工程学报，2022，42(S1)：149--161.
\end{enumerate}

\end{document}
